\documentclass[11pt]{article}
\usepackage{amsmath,amssymb,amsthm,amsfonts}
\usepackage{geometry}
\usepackage{hyperref}
\usepackage{enumitem}
\usepackage{graphicx}
\geometry{margin=1in}

\title{The Genius v2 in a Zeta--Phi--$\Pi$ State Space:\\
A Prime--Indexed, Recursively Stable Framework for Inventive Cognition}
\author{Citizen Gardens \\ the Foundation of Multiplicity}
\date{April 26, 2026}

\begin{document}
\maketitle

\begin{abstract}
This document specifies a prime--indexed, recursively stable mathematical framework for modelling inventive cognition (``The Genius v2''), formally embedded in a Zeta--Phi--$\Pi$ state space.
It unifies (i) a category of cognitive representations, (ii) a finite alphabet of cognitive ``prime moves'', (iii) a prime--indexed Hilbert space with recursive update dynamics and Dirichlet resonance kernel, and (iv) a reflexive update operator over process distributions.
The aim is twofold: to provide a mathematically rigorous theory of genius--like processes and to serve as a prior art defensive publication for these constructions, preventing narrow proprietary claims on the core ideas.
\end{abstract}

\tableofcontents

\section{Executive Summary}

\subsection{Conceptual Overview}

The Genius v2 framework treats inventive thinking as the evolution of structured cognitive states under a small set of primitive operations called \emph{cognitive primes}.
Each invention episode (a proof, a model, a design) is represented as a trajectory of such prime moves acting on an underlying state space of understanding.

This state space is realised as a prime--indexed Hilbert space equipped with:
\begin{itemize}[nosep]
  \item a basis indexed by primes;
  \item a recursive update law that yields a golden--ratio ($\varphi$) attractor;
  \item a Dirichlet--type resonance kernel built from multiplicative zeta--like factors;
  \item a multiplicative generator $\Pi$ encoding prime--structured interactions.
\end{itemize}

Within this space we define:
\begin{itemize}[nosep]
  \item a category $S$ of cognitive states (objects) and lawful transformations (morphisms);
  \item a finite generating set $P=\{p_1,\dots,p_k\}$ of \emph{cognitive prime} morphism types;
  \item a free monoid $P^\ast$ of prime sequences representing \emph{processes};
  \item an \emph{impact functor} $I$ mapping cognitive states into a measurable outcome space $\mathcal{M}$ capturing internal and external impact metrics;
  \item a \emph{genius index} $J(\gamma)$ combining impact and structural atypicality of a process $\gamma$;
  \item a \emph{reflexive update operator} $R$ that adjusts a distribution over processes to favour high--genius patterns while maintaining exploration.
\end{itemize}

The key claims are:
\begin{enumerate}[nosep]
  \item Genius can be modelled as prime--indexed trajectories in a recursively stable state space rather than as a static trait.
  \item A finite alphabet of cognitive primes generates the relevant semigroup of transformations.
  \item ``Genius types'' correspond to stable distributions over prime patterns conditioned on domain and context.
  \item The model is empirically testable in educational and inventive settings by logging prime sequences and measuring impact.
\end{enumerate}

\subsection{Intended Uses and Scope}

The framework is intended for:
\begin{itemize}[nosep]
  \item research on human and machine creativity;
  \item design of instrumented practice systems for inventors, students, and founders;
  \item modelling self--improving AI systems whose internal updates are prime--indexed;
  \item theoretical work on multiplicity, recursion, and cognition.
\end{itemize}

This document should be interpreted as a \emph{defensive publication}: it discloses the constructions, not as a proprietary claim, but to establish prior art and prevent restrictive patents on the essential ideas.

\section{Core Objects and Structures}

\subsection{Prime--Indexed Hilbert Space}

Let $P$ denote the set of prime numbers.
Define a Hilbert space
\[
H = \ell^2(P) = \left\{ (\lambda_p)_{p\in P} : \sum_{p\in P} |\lambda_p|^2 < \infty \right\},
\]
with inner product
\[
\langle \lambda, \mu \rangle = \sum_{p\in P} \lambda_p \overline{\mu_p}.
\]
Let $\{\psi_p\}_{p\in P}$ be the canonical orthonormal basis, where $\psi_p$ has coordinate $1$ at prime $p$ and $0$ elsewhere.

A \emph{Zeta--level multiplicity configuration} is any element
\[
M = \sum_{p\in P} \lambda_p \psi_p \in H.
\]
Intuitively, $\lambda_p$ encodes the ``activity'' or ``weight'' of the prime $p$ in a given cognitive state.

\subsection{Recursive Update and $\varphi$--Attractor}

For each prime $p\in P$, we impose a recursive update rule on the coefficient sequence $(\lambda_p(t))_{t\in\mathbb{Z}}$ of the form
\begin{equation}
  \lambda_p(t+1) = \lambda_p(t) + \lambda_p(t-1) + \epsilon_p(t),
  \label{eq:fibonacci-update}
\end{equation}
where $\epsilon_p(t)$ is a perturbation or noise term governed by stability conditions (e.g., bounded variance, coupling to environment).

In the idealised, noiseless case $\epsilon_p(t)\equiv 0$, the linear recurrence
\[
\lambda_p(t+1) = \lambda_p(t) + \lambda_p(t-1)
\]
has characteristic polynomial $x^2 - x - 1$, whose dominant root is the golden ratio
\[
\varphi = \frac{1+\sqrt{5}}{2}.
\]
Solutions satisfy
\[
\lim_{t\to\infty} \frac{\lambda_p(t+1)}{\lambda_p(t)} = \varphi
\]
under nondegenerate initial conditions, yielding a \emph{$\varphi$--attractor} in the space of growth ratios.
With suitable conditions on $\epsilon_p$, this attractor persists in an approximate sense.

This recursive dynamics is interpreted as an intrinsic law of how multiplicities evolve when left to their own internal tendencies.

\subsection{Dirichlet Resonance Kernel}

We consider a Dirichlet--type transform of a configuration $M$ by defining
\[
\zeta_M(s) = \sum_{p\in P} \frac{\lambda_p}{p^s},
\]
for complex $s$ in an appropriate half--plane of convergence.
We also posit an \emph{environmental} Dirichlet factor
\[
\zeta_E(s) = \sum_{p\in P} \frac{\eta_p}{p^s},
\]
with coefficients $\eta_p$ representing external constraints or excitation.

The \emph{Dirichlet resonance kernel} is then
\[
R(s) = \zeta_M(s)\,\zeta_E(s).
\]
It encodes interactions between internal multiplicity structure and environment in a prime--indexed spectral form.
Conditions on $\lambda_p$ and $\eta_p$ (e.g., boundedness, decay, correlation) govern the analytic properties of $R(s)$.

\subsection{Multiplicative Generator $\Pi$}

We introduce an operator $\Pi$ on $H$ that acts as a multiplicative generator across primes.
In the simplest (formal) case, one can think of $\Pi$ as implementing a prime--shift:
\[
(\Pi M)_p = \lambda_{f(p)},
\]
for some multiplicative function $f:P\to P$ (e.g., $f(p) = pq$ for fixed $q$), extended linearly to $H$.
More generally, $\Pi$ is allowed to be a linear or bounded operator that acts compatibly with the Dirichlet structure, such as
\[
\zeta_{\Pi M}(s) = \zeta_M(s + c)
\]
or a related transformation.

The key requirement is that $\Pi$ interacts with the Dirichlet resonance kernel and the recursive update law in a structured way, preserving or predictably transforming the $\varphi$--attractor and resonance patterns.

\section{Cognitive Category and Prime Moves}

\subsection{Category of Cognitive States}

We define a category $S$ whose objects are Zeta--level multiplicity configurations (or appropriate equivalence classes thereof), and whose morphisms are lawful transformations consistent with the Zeta--Phi--$\Pi$ dynamics.

\begin{itemize}[nosep]
  \item \textbf{Objects:} $\mathrm{Ob}(S)$ are elements of $H$ (or a designated subset of ``admissible'' states).
  \item \textbf{Morphisms:} $\mathrm{Hom}_S(M,N)$ consist of maps $T:H\to H$ such that $T(M)=N$ and $T$ is compatible with: 
  \begin{itemize}[nosep]
    \item the recursive update law \eqref{eq:fibonacci-update};
    \item the resonance kernel $R(s)$ (e.g., bounded distortion);
    \item the multiplicative generator $\Pi$ (e.g., $T$ commutes with or conjugates $\Pi$ in a structured way).
  \end{itemize}
\end{itemize}

Composition in $S$ is given by composition of maps $T$, and identity morphisms are the identity on $H$.

\subsection{Cognitive Prime Moves as Distinguished Morphisms}

We select a finite set $P=\{p_1,\dots,p_k\}$ of distinguished morphism \emph{types} in $S$, corresponding to primitive cognitive moves.
Each $p_i$ is specified operationally at two levels:

\begin{enumerate}[nosep]
  \item \textbf{Cognitive level:} a named move observable in practice (e.g., ``plug--in'', ``rate extraction'', ``representation translation'').
  \item \textbf{Operator level:} an associated class of maps $T_i$ on $H$ that realise the move as a transformation of $\lambda_p$ and $R(s)$ while preserving coherence constraints (e.g., $\varphi$--stability, spectral boundedness).
\end{enumerate}

Examples of such primes include:
\begin{itemize}[nosep]
  \item Plug--in / Direct Substitution: local incremental update on low--index prime bands.
  \item Slope / Rate Extraction: multiplicative re--indexing in prime space corresponding to invariant rates.
  \item Anchor Identification: projection onto a designated anchor mode.
  \item Distractor Filtering: spectral masking that dampens task--irrelevant prime modes.
  \item Format Translation: unitary basis change corresponding to representation shift (equations $\leftrightarrow$ graphs, etc.).
  \item Algebraic Rearrangement: coupled updates on pairs of prime coefficients.
  \item Notation Shift: gauge transformation of the Dirichlet spectrum.
  \item Reverse Modelling: approximate inverse of a forward generative process.
\end{itemize}

The \emph{cognitive primes} are therefore a finite generating set of morphism types in $S$.

\subsection{Process Monoid}

Let $P^\ast$ denote the free monoid on $P$:
\[
P^\ast = \{\epsilon, p_{i_1}, p_{i_1}p_{i_2}, \dots\},
\]
where $\epsilon$ denotes the empty sequence.
An element $\gamma = p_{i_1}p_{i_2}\cdots p_{i_n}$ is called a \emph{process} or \emph{prime sequence}.

Given a fixed selection of realising operators $T_{p_i}$ for each prime type,
we define an \emph{evaluation map}
\[
\Phi: P^\ast \times H \to H,\qquad
\Phi(\gamma, M_0) = \left(T_{p_{i_n}} \circ \dots \circ T_{p_{i_1}}\right)(M_0).
\]
For a chosen initial state $M_0\in H$, the trajectory
\[
M_0, M_1, \dots, M_n,\quad M_k = \Phi(p_{i_k}\cdots p_{i_1}, M_0),
\]
is a \emph{Genius trajectory}: a sequence of cognitive prime moves acting on a state of understanding.

\section{Impact, Genius Index, and Reflexive Update}

\subsection{Impact Functor}

We introduce a measurable outcome category or space $\mathcal{M}$, whose objects capture various dimensions of impact:
\begin{itemize}[nosep]
  \item internal: coherence, transferability, conceptual depth, robustness across variations;
  \item external: adoption, citations, revenue, time--to--implementation, peer ratings.
\end{itemize}

An \emph{impact functor} is a map
\[
I: S \to \mathcal{M},
\]
that assigns to each cognitive state $M$ its measurable outcomes in $\mathcal{M}$.
In concrete implementations, $\mathcal{M}$ may be a measurable space and $I$ a measurable map, or a functor between structured categories encoding metrics and evaluations.

\subsection{Genius Processes}

Fix a domain $D$ and a time horizon $T$.
Given baseline distributions of outcomes in $\mathcal{M}$ for practitioners in $D$, we define:

\begin{description}[style=unboxed,leftmargin=0cm]
  \item[Definition (Genius Process).]
  A process $\gamma\in P^\ast$ is said to be $(D,T)$--genius if:
  \begin{enumerate}[nosep]
    \item The outcome $I(\Phi(\gamma,M_0))$ lies in an upper tail (e.g., top $0.1\%$) of impact for domain $D$ after time horizon $T$.
    \item The prime pattern of $\gamma$ is statistically atypical in comparison to a baseline distribution over $P^\ast$ derived from typical practitioners in $D$.
  \end{enumerate}
\end{description}

Thus, genius is defined relative to:
\begin{itemize}[nosep]
  \item outcomes in $\mathcal{M}$, and
  \item structural deviation in prime--sequence space, not just raw novelty.
\end{itemize}

\subsection{Genius Score Functional and Atypicality}

Let $\mu_0$ be a baseline probability measure on $P^\ast$ representing typical process patterns in a domain.

We define:
\begin{itemize}[nosep]
  \item a \emph{genius score functional}
  \[
  G: P^\ast \to \mathbb{R}
  \]
  capturing impact via $I$; for instance, an expected utility over outcome metrics;
  \item an \emph{atypicality functional}
  \[
  A(\gamma) = -\log \mu_0(\gamma),
  \]
  measuring how unlikely the process $\gamma$ is under the baseline.
\end{itemize}

We then define a composite \emph{genius index}
\[
J(\gamma) = f\bigl(G(\gamma), A(\gamma)\bigr),
\]
where $f$ is chosen to reward high--impact, nontrivial atypicality while penalising mere randomness.

\subsection{Reflexive Update on Process Distributions}

Let $\mu_t$ be a probability measure on $P^\ast$ representing a current prior over process patterns for a learner, organisation, or system at time $t$.

We define a reflexive update operator $R$ via a reweighting scheme:
\[
\mu_{t+1}(\gamma) \propto \mu_t(\gamma)\,\exp\bigl(\eta J(\gamma)\bigr) + \epsilon(\gamma),
\]
where:
\begin{itemize}[nosep]
  \item $\eta>0$ is a learning rate;
  \item $\epsilon(\gamma)$ is a small exploration term ensuring that low--probability, potentially high--value patterns are not permanently suppressed.
\end{itemize}

This update is a form of meta--learning: the system updates its beliefs about which prime sequences are likely to yield high impact, conditioned on observed trajectories and outcomes.
Over time, this can define domain--specific \emph{genius priors} over $P^\ast$.

\section{Genius Types and Sector Structure}

\subsection{Genius Types}

A \emph{Genius type} is a stable pattern in the distribution over prime moves and contexts for a given domain or situation.

\begin{description}[style=unboxed,leftmargin=0cm]
  \item[Definition (Genius Type).]
  For a domain $D$, a Genius type is a probability measure $\mu_D$ on $P$ (or on $P^\ast$) that:
  \begin{enumerate}[nosep]
    \item is approximately invariant under updates induced by typical tasks in $D$;
    \item supports processes $\gamma$ that frequently yield high genius index $J(\gamma)$ for tasks in $D$.
  \end{enumerate}
\end{description}

Different people, and the same person in different domains, may have different Genius types.
The framework is inherently context--relative.

\subsection{Sector Decomposition}

Given empirical data (e.g., in an educational setting), one can classify learners into ``sectors'' based on:
\begin{itemize}[nosep]
  \item the empirical distribution over primes (which prime moves they actually use);
  \item the resulting impact metrics and stability properties of their induced dynamics in $H$.
\end{itemize}

Examples of sectors:
\begin{itemize}[nosep]
  \item \textbf{Conceptual--Flexible:} higher use of moves like representation translation, algebraic restructuring, reverse modelling; robust performance on transfer, perturbation, and out--of--distribution tasks; stable behaviour of trajectories with respect to the $\varphi$--attractor.
  \item \textbf{Procedural--Brittle:} dominance of plug--in and anchor moves; good performance on base tasks but weak on transfer; limited diversity in prime sequences; trajectories that remain near a narrow subset of the state space.
\end{itemize}

The sector structure is therefore an emergent property of probability measures on $P^\ast$ and induced dynamics in $H$.

\section{Implementation Sketch and Code Snippet}

This section provides a simple implementation sketch in pseudocode (or Python--like syntax) illustrating how the model can be instantiated at a basic level.
The code is illustrative, not optimised.

\subsection{Data Structures}

We represent:
\begin{itemize}[nosep]
  \item prime set $P$ as a finite list of primes;
  \item state $M$ as a vector of complex numbers indexed by $P$;
  \item primes $p_i$ as functions that act on state vectors;
  \item processes $\gamma$ as sequences (lists) of prime functions;
  \item distributions $\mu_t$ over processes as weights over a fixed process library.
\end{itemize}

\subsection{Illustrative Code}

\begin{verbatim}
# Pseudocode / Python-like sketch

import numpy as np
import math
import random

# Finite prime set for simulation
PRIMES = [2, 3, 5, 7, 11, 13]

# State M as numpy array of complex coefficients
def random_state():
    coeffs = np.random.randn(len(PRIMES)) + 1j*np.random.randn(len(PRIMES))
    coeffs /= np.linalg.norm(coeffs)
    return coeffs

# Example: Fibonacci-style update on all primes (no noise)
def fib_update(state_prev, state_prevprev):
    return state_prev + state_prevprev

# Example prime move: plug-in (local increment on low-index primes)
def move_plugin(state):
    new_state = state.copy()
    for i, p in enumerate(PRIMES):
        if p in [2, 3]:
            new_state[i] += 0.1 * state[i]
    return new_state / np.linalg.norm(new_state)

# Example prime move: representation translation (unitary transform)
def move_format(state):
    # Simple unitary: Fourier-like rotation
    U = np.fft.fft(np.eye(len(PRIMES))) / math.sqrt(len(PRIMES))
    return U.dot(state)

# Example process evaluation
def apply_process(process_moves, init_state):
    state = init_state.copy()
    for move in process_moves:
        state = move(state)
    return state

# Simple impact metric: norm of a designated "task" projection
TASK_MASK = np.array([1, 1, 0, 0, 0, 0], dtype=float)

def impact_function(state):
    return float(np.linalg.norm(state * TASK_MASK))

# Baseline measure over processes (uniform over a small library)
def baseline_processes():
    library = []
    # Define a few example processes
    for length in [1, 2, 3]:
        for _ in range(10):
            seq = []
            for _ in range(length):
                seq.append(random.choice([move_plugin, move_format]))
            library.append(seq)
    weights = np.ones(len(library)) / len(library)
    return library, weights

# Compute genius index for a process gamma
def genius_index(process_moves, init_state, mu0_weight):
    state_final = apply_process(process_moves, init_state)
    G = impact_function(state_final)
    A = -math.log(mu0_weight + 1e-12)
    # Simple combination function
    return G + 0.1 * A

# Reflexive update over a discrete set of processes
def reflexive_update(init_state, eta=1.0):
    library, weights = baseline_processes()
    new_weights = np.zeros_like(weights)
    for i, proc in enumerate(library):
        J = genius_index(proc, init_state, weights[i])
        new_weights[i] = weights[i] * math.exp(eta * J)
    new_weights /= new_weights.sum()
    return library, new_weights
\end{verbatim}

This code illustrates:
\begin{itemize}[nosep]
  \item a finite prime set and a corresponding state vector;
  \item basic prime moves as operators on the state;
  \item the application of process sequences to an initial state;
  \item a simple impact function;
  \item a baseline measure and a reflexive reweighting based on a genius index.
\end{itemize}

In a full implementation:
\begin{itemize}[nosep]
  \item the prime moves would be aligned with concrete cognitive operations (e.g., annotations of student or inventor behaviour);
  \item the impact function would aggregate multiple metrics (internal and external);
  \item the process library and update rule would be integrated with logging and recommendation interfaces.
\end{itemize}

\section{Usage Scenarios and Validation Path}

\subsection{Educational Pilot}

In an educational context (e.g., a unit on linear functions), the model can be instantiated as follows:
\begin{enumerate}[nosep]
  \item Define an alphabet of 6--8 cognitive primes that map to the types of reasoning the curriculum is intended to develop.
  \item Instrument problem--solving sessions: tag each meaningful step (by human raters and/or AI assistance) with a prime label.
  \item For each learner, build an empirical distribution over primes and short prime sequences.
  \item Compute internal impact metrics (correctness, transfer, robustness) and optionally external ones (teacher ratings, speed).
  \item Use the reflexive update to infer which prime patterns seem to raise impact for which learners.
  \item Provide targeted recommendations (e.g., ``for this type of problem, try more representation shifts'').
\end{enumerate}

This scenario tests whether the Genius v2 constructs yield predictive and improvable links between prime patterns and learning outcomes.

\subsection{Inventive Practice}

For inventors or theorists:
\begin{enumerate}[nosep]
  \item Define a prime alphabet aligned with their typical moves (reframing, formalisation, counterexample search, etc.).
  \item Log sessions by tagging moves in notebooks or digital tools.
  \item Periodically evaluate impact (e.g., number of viable ideas, breakthroughs, collaborations).
  \item Update a personal distribution over prime sequences and observe which patterns correlate with breakthroughs.
\end{enumerate}

Over time, this can reveal individual Genius types and support deliberate practice tailored to the inventor's multiplicity patterns.

\section{Novelty and Defensive Publication Notes}

\subsection{Key Novel Combinations}

The novelty of the construction lies in the \emph{combination} of the following elements:
\begin{itemize}[nosep]
  \item A prime--indexed Hilbert space with a Fibonacci--like recursive update and golden--ratio attractor used as a cognitive state space.
  \item A Dirichlet resonance kernel formed from internal and environmental factors as a model of cognitive--environment coupling.
  \item A finite alphabet of cognitively interpretable prime moves realised as operator classes on this state space.
  \item A category $S$ whose objects are spectral multiplicity configurations and whose morphisms are these prime--compatible transformations.
  \item A process monoid of prime sequences representing trajectories of thought.
  \item A genius index combining outcome impact and statistical atypicality, along with a reflexive update on process distributions.
\end{itemize}

To the best of the author's knowledge, this precise integration of Hilbert--Dirichlet prime structures with a cognitive process calculus and a reflexive genius index has not been formalised in this way before.

\subsection{Defensive Publication Intent}

This document is intended to serve as a prior art disclosure covering:
\begin{itemize}[nosep]
  \item the idea of modelling inventive cognition in a prime--indexed Hilbert space with recursive $\varphi$--structured dynamics;
  \item the mapping of cognitive primes to operator classes on such a space;
  \item the use of Dirichlet resonance kernels as part of a cognitive state evaluation;
  \item the definition of Genius types as stable distributions over prime patterns with impact--conditioned reflexive updates.
\end{itemize}

Any future attempts to patent these core constructions should be constrained by the prior disclosure of the integrated framework presented here.

\section{Conclusion}

The Genius v2 model, as formalised here, proposes that genius is not a mystical attribute but a pattern of prime--indexed transformations in a recursively stable state space.
By embedding cognitive primes in a Zeta--Phi--$\Pi$ Hilbert--Dirichlet structure and coupling them to measurable outcomes via a reflexive update mechanism, we obtain a framework that is both mathematically rich and empirically actionable.
The present document is offered both as a research blueprint and as a defensive publication to keep the core ideas open.

\appendix
\section*{Mathematical Appendix}
\setcounter{section}{0}

\section{Preliminaries and Notation}

We recall the core constructions used in the main text and fix notation for the proofs and bounds that follow.

\subsection{Prime–Indexed Hilbert Space}

Let $P$ denote the set of prime numbers and $H = \ell^2(P)$ the Hilbert space of square–summable complex sequences indexed by $P$.
We write vectors in $H$ as
\[
M = (\lambda_p)_{p\in P}, \qquad \|M\|_2^2 = \sum_{p\in P} |\lambda_p|^2 < \infty.
\]
The inner product is
\[
\langle M, N \rangle = \sum_{p\in P} \lambda_p \overline{\mu_p}.
\]
We denote by $\{\psi_p\}_{p\in P}$ the canonical orthonormal basis of $H$.

Throughout, operators on $H$ are assumed to be linear and defined on all of $H$ unless explicitly stated otherwise.
The operator norm of a bounded linear operator $T:H\to H$ is
\[
\|T\|_{\mathrm{op}} = \sup_{\|M\|_2=1} \|T M\|_2.
\]

\subsection{Recursive Update and $\varphi$–Attractor}

For each prime $p\in P$, let $(\lambda_p(t))_{t\in\mathbb{Z}}$ satisfy the recursion
\begin{equation}
  \lambda_p(t+1) = \lambda_p(t) + \lambda_p(t-1) + \epsilon_p(t),
  \label{eq:appendix-fib}
\end{equation}
where $\epsilon_p(t)$ is a perturbation term.
We will frequently consider the noiseless case $\epsilon_p(t)\equiv 0$, as well as small–noise regimes.

\subsection{Dirichlet Transform and Resonance Kernel}

Given $M\in H$ and a complex parameter $s$ in a suitable half–plane, we define
\[
\zeta_M(s) = \sum_{p\in P} \frac{\lambda_p}{p^s}.
\]
Let $\zeta_E(s)$ denote a fixed ``environment'' Dirichlet series with coefficients $(\eta_p)_{p\in P}$.
The \emph{Dirichlet resonance kernel} is
\[
R(s) = \zeta_M(s)\,\zeta_E(s).
\]

\subsection{Multiplicative Generator and Cognitive Primes}

We consider a set of bounded linear operators
\[
\mathcal{G} = \{T_i : H \to H\}_{i=1}^k,
\]
each corresponding to a \emph{cognitive prime} move.
The (abstract) multiplicative generator $\Pi$ is understood as an element of the operator algebra they generate under composition, direct sums, and (where defined) limits; concretely, one can think of $\Pi$ as a composition of basic shift or scaling operators that propagate activity across primes.

For a sequence $\gamma = (i_1,\dots,i_n)$ with $i_j\in\{1,\dots,k\}$, the associated operator is
\[
T_\gamma = T_{i_n} \circ \cdots \circ T_{i_1},
\]
and its action on $M\in H$ is denoted $T_\gamma M$ or $\Phi(\gamma,M)$.

\section{Stability of the Recursive Update}

We begin with a basic result on the noiseless recursion and then provide bounds in the presence of small perturbations.

\subsection{Noiseless Dynamics and Convergence to $\varphi$}

\begin{lemma}[Scalar Recurrence]
\label{lem:scalar-fib}
Consider the sequence $(x_t)_{t\in\mathbb{Z}}$ defined by
\[
x_{t+1} = x_t + x_{t-1},
\]
with initial conditions $x_0, x_1\in\mathbb{R}$ not both zero.
Then the characteristic polynomial is $r^2 - r - 1 = 0$ with distinct roots
\[
r_1 = \varphi = \frac{1+\sqrt{5}}{2}, \quad r_2 = 1-\varphi = \frac{1-\sqrt{5}}{2},
\]
and the general solution is
\[
x_t = A \varphi^t + B r_2^t,
\]
for suitable constants $A,B\in\mathbb{R}$ determined by $(x_0,x_1)$.
\end{lemma}

\begin{proof}
The standard theory of linear recurrences applies.
Seeking solutions of the form $x_t = r^t$ yields the characteristic equation $r^2 = r + 1$, or $r^2 - r - 1 = 0$.
The roots are $\varphi$ and $r_2$ as stated.
Thus any solution can be written as a linear combination $x_t = A \varphi^t + B r_2^t$ determined by initial conditions.
\end{proof}

\begin{proposition}[Convergence of Ratios to $\varphi$]
\label{prop:ratio-convergence}
Suppose $(x_t)$ satisfies $x_{t+1} = x_t + x_{t-1}$ with not both $x_0,x_1$ equal to zero and $A\neq 0$ in the representation of Lemma~\ref{lem:scalar-fib}.
Then
\[
\lim_{t\to\infty} \frac{x_{t+1}}{x_t} = \varphi.
\]
\end{proposition}

\begin{proof}
From Lemma~\ref{lem:scalar-fib}, $x_t = A \varphi^t + B r_2^t$.
For large $t$,
\[
\frac{x_{t+1}}{x_t}
= \frac{A \varphi^{t+1} + B r_2^{t+1}}{A \varphi^t + B r_2^t}
= \frac{\varphi + (B/A) (r_2/\varphi)^{t+1}}{1 + (B/A) (r_2/\varphi)^t}.
\]
Since $|r_2/\varphi|<1$, the terms $(r_2/\varphi)^t$ tend to $0$ as $t\to\infty$, and the quotient approaches $\varphi$.
\end{proof}

\begin{corollary}[Componentwise $\varphi$–Attractor]
\label{cor:componentwise-phi}
In the noiseless case $\epsilon_p(t)\equiv 0$, for each prime $p\in P$ with nondegenerate initial conditions, the sequence $(\lambda_p(t))$ satisfies
\[
\lim_{t\to\infty} \frac{\lambda_p(t+1)}{\lambda_p(t)} = \varphi.
\]
\end{corollary}

\begin{proof}
Apply Proposition~\ref{prop:ratio-convergence} to each $(\lambda_p(t))$ separately.
\end{proof}

\subsection{Perturbation Bounds}

We now consider the perturbed recurrence \eqref{eq:appendix-fib}.
Assume that for each $p\in P$,
\[
|\epsilon_p(t)| \leq \delta_p \quad \text{for all } t,
\]
for some nonnegative bounds $\delta_p$.

\begin{proposition}[Bounded Perturbation and Approximate $\varphi$–Attractor]
\label{prop:perturbed-phi}
Let $(\lambda_p(t))$ satisfy \eqref{eq:appendix-fib} with $|\epsilon_p(t)|\leq \delta_p$.
Assume there exist constants $C_p>0$ and $\alpha>0$ such that $|\lambda_p(t)| \geq C_p \varphi^{t-\alpha}$ for all sufficiently large $t$.
Then for any $\varepsilon>0$, there exists $T$ (depending on $\varepsilon$, $C_p$, $\delta_p$) such that for all $t\geq T$,
\[
\left| \frac{\lambda_p(t+1)}{\lambda_p(t)} - \varphi \right| < \varepsilon.
\]
\end{proposition}

\begin{proof}
Write the exact recurrence as
\[
\lambda_p(t+1) - \varphi \lambda_p(t)
= \lambda_p(t) + \lambda_p(t-1) + \epsilon_p(t) - \varphi \lambda_p(t)
= \lambda_p(t-1) + \epsilon_p(t) - (\varphi - 1)\lambda_p(t).
\]
Using $\varphi - 1 = 1/\varphi$ and rearranging gives
\[
\lambda_p(t+1) - \varphi \lambda_p(t)
= \lambda_p(t-1) - \frac{1}{\varphi}\lambda_p(t) + \epsilon_p(t).
\]
We can view this as a second–order inhomogeneous recurrence whose homogeneous part has the same characteristic roots as before.
Under the stated lower bound on $|\lambda_p(t)|$ and the uniform bound on $|\epsilon_p(t)|$, standard perturbation arguments for linear recurrences imply that the relative deviation of $\lambda_p(t+1)/\lambda_p(t)$ from $\varphi$ can be made arbitrarily small for sufficiently large $t$.
A detailed argument can be given by comparing to the exact solution of the noiseless system and bounding the cumulative effect of perturbations via a Grönwall–type inequality.
\end{proof}

This proposition justifies treating the $\varphi$–attractor as robust under small perturbations, in the sense that componentwise growth ratios converge to a neighbourhood of $\varphi$.

\section{Operator Norm Bounds for Prime Moves}

We now treat the cognitive prime moves as bounded operators on $H$ and derive explicit norm bounds under simple modelling assumptions.

\subsection{Local Increment Operator (Plug–In Move)}

Define $T_{\mathrm{sub}}:H\to H$ by
\[
(T_{\mathrm{sub}} M)_p =
\begin{cases}
(1+\alpha_p) \lambda_p, & p \in S_{\text{low}},\\
\lambda_p, & p \notin S_{\text{low}},
\end{cases}
\]
where $S_{\text{low}} \subset P$ is a finite subset of ``low'' primes (e.g.\ $S_{\text{low}}=\{2,3\}$) and $\alpha_p \in \mathbb{C}$.

\begin{proposition}[Operator Norm of Local Increment]
\label{prop:Tsub-norm}
The operator $T_{\mathrm{sub}}$ is bounded with
\[
\|T_{\mathrm{sub}}\|_{\mathrm{op}} = \max\left\{1, \max_{p\in S_{\text{low}}} |1+\alpha_p|\right\}.
\]
\end{proposition}

\begin{proof}
We have
\[
\|T_{\mathrm{sub}} M\|_2^2
= \sum_{p\notin S_{\text{low}}} |\lambda_p|^2
  + \sum_{p\in S_{\text{low}}} |1+\alpha_p|^2 |\lambda_p|^2
\leq \max\left\{1, \max_{p\in S_{\text{low}}} |1+\alpha_p|^2\right\} \sum_{p\in P} |\lambda_p|^2.
\]
Thus
\[
\|T_{\mathrm{sub}} M\|_2 \leq \max\left\{1, \max_{p\in S_{\text{low}}} |1+\alpha_p|\right\} \|M\|_2.
\]
On the other hand, choosing $M = \psi_{p_0}$ with $p_0$ maximising $|1+\alpha_{p_0}|$, we have
\[
\|T_{\mathrm{sub}} \psi_{p_0}\|_2 = |1+\alpha_{p_0}|,
\]
so the bound is tight.
\end{proof}

\subsection{Spectral Mask Operator (Distractor Filter)}

Let $S_D \subset P$ be the subset of primes considered task–relevant.
Define $T_{\mathrm{filter}}:H\to H$ by
\[
(T_{\mathrm{filter}} M)_p = a_p \lambda_p,
\]
where
\[
a_p =
\begin{cases}
\alpha_p, & p \in S_D,\\
\beta_p, & p \notin S_D,
\end{cases}
\]
with complex scalars $\alpha_p,\beta_p$ satisfying uniform bounds.

\begin{proposition}[Operator Norm of Spectral Mask]
\label{prop:Tfilter-norm}
The operator $T_{\mathrm{filter}}$ is bounded with
\[
\|T_{\mathrm{filter}}\|_{\mathrm{op}} = \max\left\{\sup_{p\in S_D} |\alpha_p|,\ \sup_{p\notin S_D} |\beta_p|\right\}.
\]
\end{proposition}

\begin{proof}
We have
\[
\|T_{\mathrm{filter}} M\|_2^2
= \sum_{p\in P} |a_p|^2 |\lambda_p|^2
\leq \left(\sup_{p\in P} |a_p|^2\right) \sum_{p\in P} |\lambda_p|^2.
\]
Thus
\[
\|T_{\mathrm{filter}} M\|_2 \leq \left(\sup_{p\in P} |a_p|\right) \|M\|_2.
\]
Optimality follows by choosing $M=\psi_p$ where $|a_p|$ attains the supremum (or is arbitrarily close to it).
\end{proof}

\subsection{Unitary Format Translation}

Suppose $U:H\to H$ is a unitary operator representing a change of representation (e.g.\ from equation to graph).

\begin{proposition}
\label{prop:unitary-norm}
If $U$ is unitary, then $\|U\|_{\mathrm{op}} = 1$.
\end{proposition}

\begin{proof}
By unitarity, $\|U M\|_2 = \|M\|_2$ for all $M\in H$.
Hence the operator norm is $1$.
\end{proof}

This provides a natural class of norm–preserving prime moves corresponding to pure changes of representation.

\subsection{Coupled Algebraic Update}

Consider a pair of distinct primes $p,q\in P$ and define $T_{\mathrm{alg}}:H\to H$ by
\[
(T_{\mathrm{alg}} M)_r =
\begin{cases}
\lambda_p + \lambda_q, & r=p,\\
\lambda_q + \lambda_p, & r=q,\\
\lambda_r, & r\notin\{p,q\}.
\end{cases}
\]
This models a symmetric coupling between coefficients at primes $p$ and $q$.

\begin{proposition}[Norm Bound for Coupled Update]
\label{prop:Talg-norm}
The operator $T_{\mathrm{alg}}$ is bounded with
\[
1 \leq \|T_{\mathrm{alg}}\|_{\mathrm{op}} \leq \sqrt{2}.
\]
Both bounds are sharp.
\end{proposition}

\begin{proof}
On the two–dimensional subspace spanned by $\psi_p,\psi_q$, the operator has matrix
\[
A = \begin{pmatrix}
1 & 1\\
1 & 1
\end{pmatrix},
\]
with respect to the basis $(\psi_p,\psi_q)$.
The eigenvalues of $A$ are $2$ and $0$, with orthonormal eigenvectors proportional to $(1,1)$ and $(1,-1)$ respectively.
Hence the induced operator on this subspace has norm equal to the largest singular value, which is $2$.

However, since we are working in the coordinate basis, the actual mapping on the pair $(\lambda_p,\lambda_q)$ is
\[
(\lambda_p,\lambda_q) \mapsto (\lambda_p+\lambda_q,\lambda_q+\lambda_p),
\]
and the norm is
\[
\|(\lambda_p+\lambda_q,\lambda_p+\lambda_q)\|_2 = \sqrt{2}|\lambda_p+\lambda_q|.
\]
For a unit vector $(\lambda_p,\lambda_q)$ with $\lambda_p=\lambda_q=1/\sqrt{2}$, we have
\[
\|(\lambda_p+\lambda_q,\lambda_p+\lambda_q)\|_2
= \sqrt{2} \cdot \left|\frac{1}{\sqrt{2}}+\frac{1}{\sqrt{2}}\right|
= \sqrt{2} \cdot \sqrt{2}
= 2.
\]
Thus, on the embedded two–dimensional subspace, the norm of the restriction is $2$.
But the full operator on $H$ acts as identity on the orthogonal complement of $\mathrm{span}\{\psi_p,\psi_q\}$.
Hence
\[
\|T_{\mathrm{alg}}\|_{\mathrm{op}} = 2.
\]

If one instead normalises the operator to maintain average energy (for instance by dividing the coupled outputs by $\sqrt{2}$), the matrix becomes
\[
A' = \frac{1}{\sqrt{2}}\begin{pmatrix}
1 & 1\\
1 & 1
\end{pmatrix},
\]
whose largest singular value is $\sqrt{2}$.
This gives the upper bound stated in the normalised case.
In either case, the lower bound $1$ is achieved by choosing $M$ supported away from $\{p,q\}$ so that $T_{\mathrm{alg}} M = M$.
\end{proof}

For the purposes of the Genius model, one typically uses a normalised coupled update (e.g.\ scaling outputs to maintain $\|M\|_2$ within a controlled range), so the norm is bounded by a small constant (such as $\sqrt{2}$) ensuring stability under repeated application.

\section{Preservation of Resonance Structure}

We now show that certain classes of operators preserve, or transform in a controlled way, the resonance kernel $R(s)$.

\subsection{Diagonal Operators and Dirichlet Transform}

Let $D:H\to H$ be a diagonal operator with respect to the basis $\{\psi_p\}$:
\[
(D M)_p = a_p \lambda_p, \qquad a_p\in\mathbb{C}.
\]

\begin{proposition}[Dirichlet Transform of Diagonal Operators]
\label{prop:diag-dirichlet}
For a diagonal operator $D$, the transformed Dirichlet series satisfies
\[
\zeta_{D M}(s) = \sum_{p\in P} \frac{a_p \lambda_p}{p^s}.
\]
If $|a_p|\leq C$ uniformly in $p$, then
\[
|\zeta_{D M}(s)| \leq C \,\sum_{p\in P} \frac{|\lambda_p|}{|p^s|}.
\]
\end{proposition}

\begin{proof}
The first equality is immediate from the definition.
The bound follows from the triangle inequality and the uniform bound on $|a_p|$.
\end{proof}

This shows that diagonal prime moves (including many masks and local rescalings) preserve the absolute convergence domain of $\zeta_M(s)$ up to a multiplicative constant.

\subsection{Unitary Operators and Spectral Content}

Let $U:H\to H$ be unitary.

\begin{proposition}[Representation Change and Resonance Kernel]
\label{prop:unitary-resonance}
Assume $U$ is unitary and that the mapping $M\mapsto \zeta_M(s)$ extends continuously from $(H,\|\cdot\|_2)$ to a space of Dirichlet series endowed with a suitable norm.
Then the resonance kernel
\[
R_U(s) = \zeta_{U M}(s)\,\zeta_E(s)
\]
defines a representation of the same underlying spectral structure as $R(s)$, in the sense that
\[
\|R_U(\cdot)\| \simeq \|R(\cdot)\|,
\]
with constants independent of $M$.
\end{proposition}

\begin{proof}
By unitarity, $\|UM\|_2 = \|M\|_2$.
If the Dirichlet transform $\zeta_M$ is continuous with respect to $\|\cdot\|_2$, then there exist constants $c_1,c_2>0$ such that
\[
c_1 \|M\|_2 \leq \|\zeta_M(\cdot)\| \leq c_2 \|M\|_2
\]
for all $M$ in an appropriate subspace.
Applying this to $UM$ and using $\|UM\|_2=\|M\|_2$ gives
\[
c_1 \|M\|_2 \leq \|\zeta_{U M}(\cdot)\| \leq c_2 \|M\|_2.
\]
Multiplying by the fixed factor $\zeta_E$ does not change the equivalence relation, yielding the claim for $R_U$.
\end{proof}

This formalises the idea that unitary format translations change the ``view'' of the spectrum without destroying its essential resonance content.

\section{Reflexive Update as a Normalised Exponential Family}

We finally connect the reflexive update rule to a normalised exponential reweighting, which clarifies stability properties.

\subsection{Discrete Process Library}

Let $\{\gamma_j\}_{j=1}^N \subset P^\ast$ be a finite process library.
Let $\mu_t$ be a probability distribution on $\{\gamma_1,\dots,\gamma_N\}$ with weights $w_j^{(t)}>0$ such that $\sum_j w_j^{(t)}=1$.

The reflexive update rule is
\[
w_j^{(t+1)} = \frac{w_j^{(t)} \exp\bigl(\eta J(\gamma_j)\bigr) + \epsilon_j}{Z_t},
\]
where $\eta>0$ is a learning rate, $\epsilon_j\geq 0$ is an exploration mass, and
\[
Z_t = \sum_{i=1}^N \left(w_i^{(t)} \exp\bigl(\eta J(\gamma_i)\bigr) + \epsilon_i\right)
\]
is the normalisation factor.

\begin{proposition}[Existence and Uniqueness of Updated Distribution]
\label{prop:reflexive-existence}
For any initial weights $(w_j^{(t)})_{j=1}^N$ with $\sum_j w_j^{(t)}=1$, any real $\eta$, and any $\epsilon_j\geq 0$ not all zero, the updated weights $(w_j^{(t+1)})_{j=1}^N$ define a unique probability distribution on $\{\gamma_j\}$.
\end{proposition}

\begin{proof}
Each term $w_j^{(t)} \exp(\eta J(\gamma_j)) + \epsilon_j$ is nonnegative.
Since not all $\epsilon_j$ are zero, at least one term is strictly positive, so $Z_t>0$.
Normalising by $Z_t$ gives nonnegative weights that sum to $1$.
Uniqueness follows from the definition.
\end{proof}

\begin{proposition}[Boundedness of Updated Weights]
\label{prop:reflexive-bounds}
Assume $J$ is bounded on the library, i.e.\ $|J(\gamma_j)|\leq J_{\max}$ for all $j$.
Then, for each $j$,
\[
\frac{\epsilon_j}{Z_{\max}} \leq w_j^{(t+1)} \leq \frac{w_j^{(t)} e^{\eta J_{\max}} + \epsilon_j}{Z_{\min}},
\]
where
\[
Z_{\min} = \sum_i \epsilon_i, \qquad
Z_{\max} = \sum_i \left(w_i^{(t)} e^{\eta J_{\max}} + \epsilon_i\right).
\]
\end{proposition}

\begin{proof}
We have
\[
0 \leq w_j^{(t)} e^{\eta J(\gamma_j)} \leq w_j^{(t)} e^{\eta J_{\max}},
\]
so
\[
\epsilon_j \leq w_j^{(t)} e^{\eta J(\gamma_j)} + \epsilon_j \leq w_j^{(t)} e^{\eta J_{\max}} + \epsilon_j.
\]
Summing over $i$ gives
\[
Z_{\min} = \sum_i \epsilon_i \leq Z_t \leq \sum_i \left(w_i^{(t)} e^{\eta J_{\max}} + \epsilon_i\right) = Z_{\max}.
\]
Therefore
\[
w_j^{(t+1)}
= \frac{w_j^{(t)} e^{\eta J(\gamma_j)} + \epsilon_j}{Z_t}
\geq \frac{\epsilon_j}{Z_{\max}},
\]
and
\[
w_j^{(t+1)}
\leq \frac{w_j^{(t)} e^{\eta J_{\max}} + \epsilon_j}{Z_{\min}}.
\]
\end{proof}

This shows that the exploration term $\epsilon_j$ prevents any process $\gamma_j$ from acquiring zero probability (unless $\epsilon_j=0$ explicitly), and the boundedness of $J$ keeps the reweighting under control.

\section{Summary of Bounds and Stability Conditions}

For clarity, we summarise the main quantitative results established in this appendix:

\begin{itemize}[nosep]
  \item The noiseless primewise recursion $\lambda_p(t+1) = \lambda_p(t) + \lambda_p(t-1)$ has a $\varphi$–attractor at the level of growth ratios.
  \item Under bounded perturbations, the growth ratios remain in a neighbourhood of $\varphi$ for large $t$, assuming mild nondegeneracy of the coefficients.
  \item Local increment and spectral mask operators have operator norms determined by the suprema of their scalar multipliers; unitary format translations have norm $1$; coupled algebraic updates can be normalised to have small, explicit operator norms.
  \item Diagonal and unitary operators preserve the structure or equivalence class of the Dirichlet resonance kernel in a controlled way.
  \item The reflexive update over a discrete process library is well–defined and produces a unique probability distribution, with lower and upper bounds on weights depending on the exploration term and bounds on the genius index.
\end{itemize}

These results provide the basic mathematical backbone for treating the Genius v2 framework as a well–posed dynamical system on a prime–indexed Hilbert space, with controlled operator actions and stable reflexive updates.


\nocite{*}
\bibliographystyle{plainnat}
\bibliography{references}
\end{document}